\chapter{Results and Discussion}
\label{chap:results}

This chapter presents the results of all experiments conducted in this research, following the zigzag structure with the methodology described in Chapter~\ref{chap:methodology}. Section~\ref{sec:res_baseline} reports the baseline retrieval performance. Section~\ref{sec:res_error} presents the error analysis findings that guided technique selection. Section~\ref{sec:res_query2doc} reports the initial Query2Doc results with both dense and sparse retrieval. Section~\ref{sec:res_model_comparison} presents the comprehensive model comparison across ten LLMs. Finally, Section~\ref{sec:res_key_findings} synthesises the key findings and their implications.

%======================================================================
\section{Baseline Results and Comparison}
\label{sec:res_baseline}
%======================================================================

The two retrieval baselines described in Section~\ref{sec:meth_baseline} were evaluated on the full MIRACL Arabic development set (2,896 queries). Table~\ref{tab:baseline_results} presents the results.

\begin{table}[htbp]
    \centering
    \caption{Baseline retrieval results on MIRACL Arabic dev set (2,896 queries). Bold values indicate the higher score for each metric.}
    \label{tab:baseline_results}
    \begin{tabular}{@{}lcccc@{}}
        \toprule
        \textbf{Retriever} & \textbf{NDCG@10} & \textbf{Recall@10} & \textbf{Recall@100} & \textbf{MRR} \\
        \midrule
        mDPR (Dense) & \textbf{0.4993} & \textbf{0.6156} & 0.8407 & \textbf{0.5328} \\
        BM25S (Sparse) & 0.4621 & 0.5964 & \textbf{0.8577} & 0.4836 \\
        \bottomrule
    \end{tabular}
\end{table}

The mDPR dense baseline reproduced the official MIRACL results with less than 0.1\% difference (NDCG@10 = 0.4993 vs. published 0.499), confirming the correctness of the implementation. The BM25S sparse baseline achieved 96\% of the official Pyserini BM25 performance (NDCG@10 = 0.4621 vs. target 0.481), with the 4\% difference attributable to the use of BM25S rather than Pyserini's Java-based implementation, as discussed in Section~\ref{sec:meth_bm25_baseline}.

\subsection{Complementary Strengths}
\label{sec:res_baseline_comparison}

The two baselines exhibited complementary performance characteristics. mDPR achieved higher NDCG@10 (+8.1\%) and MRR (+10.2\%), indicating superior ranking quality---relevant passages were placed higher in the result list. BM25S achieved higher Recall@100 (+2.0\%), indicating broader retrieval coverage---more relevant passages were retrieved overall, though ranked less effectively.

This complementary behaviour aligns with the theoretical properties of dense and sparse retrieval discussed in Section~\ref{sec:ir_methods}: dense retrieval captures semantic similarity and ranks well, while sparse retrieval excels at broad lexical matching. This finding confirmed the decision to test query enhancement techniques on both retrievers independently.

%======================================================================
\section{Error Analysis Findings}
\label{sec:res_error}
%======================================================================

The error analysis methodology described in Section~\ref{sec:meth_error} was applied to the dense baseline (Experiment~001). The findings guided the selection of the first query enhancement technique.

\subsection{Overall Failure Rate}
\label{sec:res_error_rate}

Table~\ref{tab:error_segmentation} presents the performance segmentation of all 2,896 queries.

\begin{table}[htbp]
    \centering
    \caption{Performance segmentation of baseline queries by NDCG@10 score.}
    \label{tab:error_segmentation}
    \begin{tabular}{@{}lccc@{}}
        \toprule
        \textbf{Category} & \textbf{Count} & \textbf{\% of Total} & \textbf{Avg NDCG@10} \\
        \midrule
        Failed (NDCG@10 $< 0.3$) & 1,130 & 39.0\% & 0.158 \\
        Mediocre ($0.3 \leq$ NDCG@10 $< 0.7$) & 1,618 & 55.9\% & 0.472 \\
        Successful (NDCG@10 $\geq 0.7$) & 148 & 5.1\% & 0.762 \\
        \bottomrule
    \end{tabular}
\end{table}

A striking finding was that 39\% of queries (1,130 out of 2,896) were classified as failed, with 192 queries achieving zero recall in the top-10 (no relevant passages retrieved). Only 5.1\% of queries were classified as highly successful. This high failure rate established a clear opportunity for query enhancement to make a measurable impact.

\subsection{Short Query Performance Gap}
\label{sec:res_error_length}

The analysis of retrieval performance by query length revealed a significant performance gap, as shown in Table~\ref{tab:query_length}.

\begin{table}[htbp]
    \centering
    \caption{Retrieval performance by query length bucket.}
    \label{tab:query_length}
    \begin{tabular}{@{}lccc@{}}
        \toprule
        \textbf{Length Bucket} & \textbf{Avg NDCG@10} & \textbf{Count} & \textbf{\% of Total} \\
        \midrule
        Short (1--3 tokens) & 0.240 & 147 & 5.1\% \\
        Medium (4--8 tokens) & 0.367 & 2,495 & 86.2\% \\
        Long (9+ tokens) & 0.406 & 254 & 8.8\% \\
        \bottomrule
    \end{tabular}
\end{table}

Short queries achieved only 59\% of the performance of long queries (NDCG@10 = 0.240 vs. 0.406), representing a 41\% performance gap. The Pearson correlation between query length and NDCG@10 was $r = 0.125$ ($p < 0.001$), indicating a weak but statistically significant positive relationship. This finding directly motivated the selection of query expansion as the first enhancement technique: by generating pseudo-documents that add vocabulary and context, the effective information content of short queries would be increased.

\subsection{Retrieval Coverage Analysis}
\label{sec:res_error_coverage}

The coverage analysis revealed an important distinction between retrieval and ranking:

\begin{table}[htbp]
    \centering
    \caption{Proportion of queries with at least one relevant passage at each retrieval depth.}
    \label{tab:coverage}
    \begin{tabular}{@{}lc@{}}
        \toprule
        \textbf{Retrieval Depth} & \textbf{Coverage} \\
        \midrule
        Top-10 & 93.4\% \\
        Top-20 & 96.7\% \\
        Top-50 & 98.8\% \\
        Top-100 & 99.4\% \\
        \bottomrule
    \end{tabular}
\end{table}

Retrieval coverage was excellent at depth 100 (99.4\%), but dropped to 93.4\% at depth 10. This 6.0 percentage point gap indicated that the retriever was \emph{finding} relevant passages but not \emph{ranking} them highly enough. Query enhancement could address this by enriching the query representation to produce better-calibrated similarity scores.

\subsection{Technique Selection Rationale}
\label{sec:res_error_rationale}

Manual inspection of the 20 worst-performing queries revealed recurring failure patterns: vocabulary mismatch (e.g., the Arabic term ``\<آزوت>'' for nitrogen versus the more common ``\<نيتروجين>''), named entity variations (e.g., ``\<إبن الهيثم>'' vs. ``\<ابن الهيثم>''), and diacritics sensitivity (e.g., ``\<المَثَانةُ>'' vs. ``\<المثانة>''). Based on these findings, Query2Doc \cite{wang_2023_query2doc} was selected as the first technique because its pseudo-document generation mechanism directly addresses information poverty in short queries and vocabulary mismatch---the two dominant failure patterns identified. The complete rationale for this selection is documented in Section~\ref{sec:meth_query2doc}.

%======================================================================
\section{Query2Doc Results}
\label{sec:res_query2doc}
%======================================================================

This section presents the results of the initial Query2Doc experiments using Qwen~2.5~3B~Instruct, evaluated on both dense and sparse retrieval.

\subsection{Dense Retrieval (Experiment 003)}
\label{sec:res_q2d_dense}

Table~\ref{tab:q2d_dense} presents the results of Query2Doc with dense retrieval.

\begin{table}[htbp]
    \centering
    \caption{Query2Doc results with dense retrieval (mDPR). Improvement is computed relative to the dense baseline (Experiment~001).}
    \label{tab:q2d_dense}
    \begin{tabular}{@{}lccc@{}}
        \toprule
        \textbf{Metric} & \textbf{Baseline} & \textbf{Query2Doc} & \textbf{Improvement} \\
        \midrule
        NDCG@10 & 0.4993 & 0.5435 & +8.9\% \\
        Recall@10 & 0.6156 & 0.6608 & +7.3\% \\
        Recall@100 & 0.8407 & 0.8594 & +2.2\% \\
        MRR & 0.5328 & 0.5742 & +7.8\% \\
        \bottomrule
    \end{tabular}
\end{table}

Query2Doc produced a substantial improvement of +8.9\% in NDCG@10 over the dense baseline, representing the largest single-technique improvement achieved in the initial experiments. All metrics improved, with NDCG@10 and Recall@10 showing the largest gains. The smaller Recall@100 improvement (+2.2\%) indicates that Query2Doc primarily improved \emph{ranking} (moving relevant documents higher) rather than \emph{discovery} (finding new relevant documents).

The query expansion statistics showed consistent enhancement: the median query length increased from 27 characters (original) to 250 characters (enhanced), representing an 8.45$\times$ expansion ratio. The LLM-generated pseudo-documents provided additional vocabulary and semantic context that improved the mDPR encoder's ability to match queries with relevant passages.

The total runtime was approximately 40 minutes for 2,896 queries on a T4 GPU, achieved through the batch processing optimisations described in Section~\ref{sec:meth_q2d_batch}.

\subsection{BM25 Retrieval (Experiment 004)}
\label{sec:res_q2d_bm25}

Table~\ref{tab:q2d_bm25} presents the results of Query2Doc with BM25S sparse retrieval.

\begin{table}[htbp]
    \centering
    \caption{Query2Doc results with BM25 retrieval (BM25S). Negative values indicate performance degradation.}
    \label{tab:q2d_bm25}
    \begin{tabular}{@{}lccc@{}}
        \toprule
        \textbf{Metric} & \textbf{Baseline} & \textbf{Query2Doc} & \textbf{Change} \\
        \midrule
        NDCG@10 & 0.4621 & 0.4090 & $-$11.5\% \\
        Recall@10 & 0.5964 & 0.5384 & $-$9.7\% \\
        Recall@100 & 0.8577 & 0.8155 & $-$4.9\% \\
        MRR & 0.4836 & 0.4342 & $-$10.2\% \\
        \bottomrule
    \end{tabular}
\end{table}

In stark contrast to the dense results, Query2Doc \emph{degraded} BM25 performance across all metrics, with NDCG@10 declining by 11.5\%. This result, while negative, provided a valuable insight into the interaction between query enhancement techniques and retrieval paradigms.

\subsubsection{Term Dilution Analysis}
\label{sec:res_term_dilution}

The performance degradation was attributed to \textbf{term dilution}: the BM25 scoring function (Equation~\ref{eq:bm25}) is sensitive to query length, as additional terms distribute the term-frequency weight across a larger vocabulary. When a short query such as ``\<ما هي عاصمة السودان>'' (``What is the capital of Sudan,'' 4 tokens) was expanded to 200+ tokens, the importance of the original keywords ``\<عاصمة>'' (capital) and ``\<السودان>'' (Sudan) was diluted in the BM25 scoring.

As noted in Section~\ref{sec:meth_q2d_bm25}, this experiment did not implement the query repetition strategy recommended by Wang et al. \cite{wang_2023_query2doc}, where the original query is repeated $n=5$ times before concatenation. This omission was the primary cause of the BM25 degradation. The term dilution phenomenon and its resolution are further examined in the model comparison experiments (Section~\ref{sec:res_mc_bm25}), where models with Arabic-specific vocabulary were observed to partially overcome this limitation.

\subsection{Comparison with the Original Query2Doc Paper}
\label{sec:res_q2d_comparison}

Table~\ref{tab:q2d_paper_comparison} compares the results of this implementation with those reported in the original Query2Doc paper \cite{wang_2023_query2doc}.

\begin{table}[htbp]
    \centering
    \caption{Comparison between the original Query2Doc results (English, MS-MARCO) and this implementation (Arabic, MIRACL).}
    \label{tab:q2d_paper_comparison}
    \begin{tabular}{@{}lcc@{}}
        \toprule
        \textbf{Configuration} & \textbf{Original Paper} & \textbf{This Work} \\
        \midrule
        Language & English & Arabic \\
        Dataset & MS-MARCO & MIRACL Arabic \\
        LLM & GPT-3 (175B) & Qwen 2.5 3B \\
        Prompting & Few-shot (4 examples) & Zero-shot \\
        Dense NDCG@10 improvement & +2--5\% & +8.9\% \\
        Sparse NDCG@10 improvement & +3--7\% & $-$11.5\%\textsuperscript{*} \\
        \bottomrule
    \end{tabular}
    \\[0.5em]
    \raggedright\small\textsuperscript{*}Without query repetition; the original paper uses $n=5$ repetition for sparse retrieval.
\end{table}

The dense retrieval improvement (+8.9\%) exceeded the original paper's reported range (+2--5\%), despite using a model that is 58$\times$ smaller and zero-shot rather than few-shot prompting. Two hypotheses were considered for this result: (1) Arabic queries, with their morphological richness (Section~\ref{sec:arabic_challenges}), may benefit disproportionately from the vocabulary expansion provided by pseudo-documents; and (2) the mDPR baseline, being a weaker retriever not fine-tuned on MIRACL, may have more room for improvement than the strong baselines used in the original paper.

%======================================================================
\section{Model Comparison Results}
\label{sec:res_model_comparison}
%======================================================================

Following the methodology described in Section~\ref{sec:meth_model_comparison}, ten open-source LLMs were evaluated using the identical Query2Doc pipeline. This section presents the results for both dense and sparse retrieval, including the analysis of dropped models.

\subsection{Dense Retrieval Leaderboard}
\label{sec:res_mc_dense}

Table~\ref{tab:dense_leaderboard} presents the complete dense retrieval results, ranked by NDCG@10.

\begin{table}[htbp]
    \centering
    \caption{Dense retrieval (mDPR) results for all models, ranked by NDCG@10. The improvement column is relative to the no-enhancement baseline (Experiment~001).}
    \label{tab:dense_leaderboard}
    \begin{tabular}{@{}clccccc@{}}
        \toprule
        \textbf{Rank} & \textbf{Model} & \textbf{Params} & \textbf{NDCG@10} & \textbf{Recall@10} & \textbf{MRR} & \textbf{$\Delta$NDCG} \\
        \midrule
        --- & Baseline (no QE) & --- & 0.4993 & 0.6156 & 0.5328 & --- \\
        \midrule
        1 & Aya Expanse 8B & 8B & \textbf{0.6164} & \textbf{0.7256} & \textbf{0.6493} & +23.5\% \\
        2 & Jais-2 8B & 8B & 0.6018 & 0.7161 & 0.6356 & +20.5\% \\
        3 & Qwen3-8B & 8B & 0.5958 & 0.7119 & 0.6278 & +19.3\% \\
        4 & Qwen 2.5 7B & 7B & 0.5813 & 0.6952 & 0.6134 & +16.4\% \\
        5 & Qwen3-4B & 4B & 0.5691 & 0.6824 & 0.6015 & +14.0\% \\
        6 & Qwen 2.5 3B & 3B & 0.5435 & 0.6608 & 0.5742 & +8.9\% \\
        7 & Gemma 3 4B & 4B & 0.5391 & 0.6443 & 0.5761 & +8.0\% \\
        8 & Falcon-H1 3B & 3B & 0.5359 & 0.6484 & 0.5681 & +7.3\% \\
        9 & SILMA 2B & 2B & 0.5177 & 0.6289 & 0.5508 & +3.7\% \\
        \midrule
        \multicolumn{7}{@{}l}{\textit{Dropped models (see Section~\ref{sec:res_dropped}):}} \\
        --- & ALLaM 7B & 7B & 0.2550 & 0.3335 & 0.2708 & $-$48.9\% \\
        --- & GPT-OSS 20B & 20B & \multicolumn{4}{c}{Dropped before full evaluation} \\
        \bottomrule
    \end{tabular}
\end{table}

All nine successfully evaluated models improved over the dense baseline, with improvements ranging from +3.7\% (SILMA 2B) to +23.5\% (Aya Expanse 8B). The top three models---Aya Expanse 8B, Jais-2 8B, and Qwen3-8B---all achieved improvements exceeding +19\%, demonstrating that Query2Doc with appropriately selected LLMs can substantially enhance Arabic dense retrieval.

% TODO: Add bar chart figure
\begin{figure}[htbp]
    \centering
    \fbox{\parbox{0.85\textwidth}{\centering\vspace{2em}
    \textbf{[Figure placeholder: Dense Retrieval NDCG@10 Bar Chart]}\\[0.5em]
    \small Horizontal bar chart showing all models ranked by NDCG@10, with baseline indicated by a dashed vertical line.
    \vspace{2em}}}
    \caption{Dense retrieval NDCG@10 for all models. The dashed line indicates the no-enhancement baseline (0.4993). ALLaM is shown separately as a dropped model.}
    \label{fig:dense_bar_chart}
\end{figure}

\subsection{BM25 Retrieval Results}
\label{sec:res_mc_bm25}

Table~\ref{tab:bm25_leaderboard} presents the BM25 retrieval results.

\begin{table}[htbp]
    \centering
    \caption{BM25 retrieval (BM25S) results for all models, ranked by NDCG@10. The improvement column is relative to the no-enhancement baseline (Experiment~002). Negative values indicate degradation.}
    \label{tab:bm25_leaderboard}
    \begin{tabular}{@{}clccccc@{}}
        \toprule
        \textbf{Rank} & \textbf{Model} & \textbf{Params} & \textbf{NDCG@10} & \textbf{Recall@10} & \textbf{MRR} & \textbf{$\Delta$NDCG} \\
        \midrule
        --- & Baseline (no QE) & --- & 0.4621 & 0.5964 & 0.4836 & --- \\
        \midrule
        1 & Jais-2 8B & 8B & \textbf{0.5122} & \textbf{0.6448} & \textbf{0.5397} & +10.8\% \\
        2 & Aya Expanse 8B & 8B & 0.5046 & 0.6284 & 0.5377 & +9.2\% \\
        3 & Qwen 2.5 7B & 7B & 0.4682 & 0.6040 & 0.4905 & +1.3\% \\
        \midrule
        4 & Qwen3-8B & 8B & 0.4459 & 0.5806 & 0.4702 & $-$3.5\% \\
        5 & SILMA 2B & 2B & 0.4277 & 0.5550 & 0.4485 & $-$7.4\% \\
        6 & Qwen3-4B & 4B & 0.4145 & 0.5403 & 0.4415 & $-$10.3\% \\
        7 & Qwen 2.5 3B & 3B & 0.4090 & 0.5384 & 0.4342 & $-$11.5\% \\
        8 & Falcon-H1 3B & 3B & 0.4038 & 0.5311 & 0.4274 & $-$12.6\% \\
        9 & Gemma 3 4B & 4B & 0.3447 & 0.4532 & 0.3718 & $-$25.4\% \\
        \bottomrule
    \end{tabular}
\end{table}

The BM25 results were markedly different from the dense results. Only three models improved over the BM25 baseline: Jais-2 8B (+10.8\%), Aya Expanse 8B (+9.2\%), and Qwen~2.5 7B (+1.3\%). The remaining six models degraded BM25 performance, with Gemma~3~4B showing the largest decline ($-$25.4\%).

The fact that two models---Jais-2 and Aya---achieved substantial BM25 improvements \emph{without} implementing the query repetition strategy (Equation~\ref{eq:q2d_bm25}) is a notable finding. Jais-2's success with BM25 is hypothesised to stem from its Arabic-centric vocabulary of 150,000 tokens (described in Section~\ref{sec:jais2}), which produces pseudo-documents with vocabulary that closely matches the BM25 index terms. Similarly, Aya Expanse's purpose-built multilingual training (101 languages with explicit Arabic optimisation, Section~\ref{sec:aya}) appears to generate lexically precise expansions.

\subsection{Dropped Models Analysis}
\label{sec:res_dropped}

Two models were dropped from the final comparison due to severe performance issues. Following Dr.~Tahani's guidance that even negative results constitute valuable research contributions, both cases are analysed in detail.

\subsubsection{ALLaM-7B (Experiment 008)}
\label{sec:res_allam}

ALLaM-7B (described in Section~\ref{sec:allam}) produced a catastrophic $-$48.9\% decline in NDCG@10 (0.2550 vs. baseline 0.4993), performing \emph{worse} than using no query enhancement at all. Investigation revealed a sentencepiece tokeniser artefact: the Unicode character ``\texttt{\char"2581}'' (used internally to mark word boundaries) leaked into the decoded pseudo-documents. These special characters contaminated the query embeddings, causing the mDPR encoder to produce degraded representations.

This failure was attributed to ALLaM's preview/alpha release status at the time of testing. The model had been accepted at ICLR~2025 \cite{bari2025allam} but was not yet in stable release, and the tokeniser behaviour had not been fully validated for generation tasks.

\subsubsection{GPT-OSS-20B (Experiment 009)}
\label{sec:res_gptoss}

GPT-OSS-20B (described in Section~\ref{sec:gptoss}) was dropped during the sanity check phase due to three compounding issues:

\begin{enumerate}
    \item \textbf{Impractical inference speed:} The Mixture-of-Experts architecture (32 experts, top-4 routing) was approximately 70$\times$ slower than dense models of comparable active parameter count when run through 4-bit quantisation. A full run of 2,896 queries would have required an estimated 40+ hours.
    \item \textbf{English reasoning leakage:} The Harmony chat format caused the model to produce English-language reasoning before Arabic output, requiring a forced-final-channel prefix to suppress. Even with this workaround, output quality was inconsistent.
    \item \textbf{Factual hallucination:} In the 5-query sanity check, 3 out of 5 generated pseudo-documents contained severe factual errors despite producing 100\% Arabic output. Hallucinated content would introduce misleading vocabulary into the enhanced queries.
\end{enumerate}

These issues demonstrated that the model's English-dominant training data renders it unreliable for Arabic query enhancement, regardless of its large parameter count.

%======================================================================
\section{Key Findings and Analysis}
\label{sec:res_key_findings}
%======================================================================

This section synthesises the results across all experiments to identify the key findings and their implications.

\subsection{Model Size Correlates with Dense Retrieval Performance}
\label{sec:res_finding_size}

A clear positive correlation was observed between model parameter count and dense retrieval improvement. Figure~\ref{fig:size_vs_performance} illustrates this relationship.

% TODO: Add scatter plot figure
\begin{figure}[htbp]
    \centering
    \fbox{\parbox{0.85\textwidth}{\centering\vspace{2em}
    \textbf{[Figure placeholder: Model Size vs. NDCG@10 Improvement]}\\[0.5em]
    \small Scatter plot with model size (B) on x-axis and NDCG@10 improvement (\%) on y-axis. Points labelled by model name. Trend line showing positive correlation. ALLaM excluded as outlier.
    \vspace{2em}}}
    \caption{Relationship between model parameter count and dense retrieval NDCG@10 improvement. ALLaM-7B is excluded as an outlier due to tokeniser artefacts.}
    \label{fig:size_vs_performance}
\end{figure}

Grouping models by size tier reveals a consistent pattern:

\begin{itemize}
    \item \textbf{2--3B models:} +3.7\% to +8.9\% NDCG@10 improvement (SILMA, Qwen~2.5~3B, Falcon-H1)
    \item \textbf{4B models:} +8.0\% to +14.0\% (Gemma~3, Qwen3-4B)
    \item \textbf{7--8B models:} +16.4\% to +23.5\% (Qwen~2.5~7B, Qwen3-8B, Jais-2, Aya)
\end{itemize}

This suggests that larger models generate higher-quality pseudo-documents with richer vocabulary and more accurate factual content, leading to better-calibrated query representations in the dense embedding space.

\subsection{Generational Improvement Within Model Families}
\label{sec:res_finding_generation}

The Qwen model family, with four models tested across two generations, provided a controlled comparison of generational improvements:

\begin{table}[htbp]
    \centering
    \caption{Generational comparison within the Qwen model family.}
    \label{tab:qwen_generations}
    \begin{tabular}{@{}lccccc@{}}
        \toprule
        \textbf{Comparison} & \textbf{Qwen 2.5} & \textbf{Qwen3} & \textbf{$\Delta$} & \textbf{Training Data} \\
        \midrule
        3--4B (Dense NDCG@10) & 0.5435 & 0.5691 & +4.7\% & 18T $\rightarrow$ 36T \\
        7--8B (Dense NDCG@10) & 0.5813 & 0.5958 & +2.5\% & 18T $\rightarrow$ 36T \\
        \bottomrule
    \end{tabular}
\end{table}

At both parameter scales, the newer Qwen3 generation outperformed Qwen~2.5, despite identical prompting and evaluation conditions. The improvement is attributed to Qwen3's doubled training data volume (36T vs. 18T tokens), expanded language coverage (119 vs. 29 languages including 8 Arabic dialects), and wider attention architecture (described in Sections~\ref{sec:qwen3_4b} and \ref{sec:qwen3_8b}).

\subsection{Arabic Specialisation vs. Multilingual Training}
\label{sec:res_finding_arabic}

An important finding was that Arabic-specialised models did not consistently outperform multilingual models. Falcon-H1-3B, despite holding the highest Arabic OALL benchmark score ($\sim$62\%) at the 3B scale, was outperformed on dense NDCG@10 by both Qwen~2.5~3B (a multilingual model with a lower Arabic benchmark score) and Qwen3-4B. Similarly, Jais-2-8B, while achieving the best BM25 results, was surpassed on dense retrieval by Aya~Expanse~8B, a model designed for 101 languages without Arabic-specific architecture.

This suggests that \textbf{Arabic NLP benchmark scores do not directly predict query expansion quality}. The skills measured by Arabic benchmarks (reading comprehension, question answering, reasoning) differ from the ability to generate vocabulary-rich, lexically diverse pseudo-documents that improve retrieval. Training data volume and diversity appear to be more important factors for this specific task.

The notable exception is Jais-2's BM25 performance: its custom 150,000-token Arabic vocabulary (Section~\ref{sec:jais2}) produces expansions with term distributions that align closely with BM25 index terms, making it uniquely effective for sparse retrieval.

\subsection{Dense vs. BM25 Behaviour with Query Enhancement}
\label{sec:res_finding_retriever}

The divergent response of dense and sparse retrieval to query enhancement is one of the most practically significant findings of this research. Table~\ref{tab:retriever_divergence} summarises the pattern.

\begin{table}[htbp]
    \centering
    \caption{Number of models improving over baseline, by retriever type.}
    \label{tab:retriever_divergence}
    \begin{tabular}{@{}lcc@{}}
        \toprule
        \textbf{Outcome} & \textbf{Dense (mDPR)} & \textbf{BM25S} \\
        \midrule
        Models improving & 9 / 9 & 3 / 9 \\
        Models degrading & 0 / 9 & 6 / 9 \\
        Avg improvement & +12.3\% & $-$5.5\% \\
        Best improvement & +23.5\% & +10.8\% \\
        Worst degradation & +3.7\% & $-$25.4\% \\
        \bottomrule
    \end{tabular}
\end{table}

Dense retrieval benefited universally from query enhancement because the mDPR encoder operates in a continuous semantic space where additional context improves the query embedding representation. BM25, conversely, relies on exact term matching, where additional terms can dilute the weights of the original query keywords (Section~\ref{sec:res_term_dilution}).

This finding has direct implications for RAG system design: query enhancement techniques must be adapted to the retrieval paradigm, and a strategy that works well for dense retrieval should not be assumed to transfer to sparse retrieval without modification.

\subsection{Best Model Recommendations}
\label{sec:res_finding_best}

Based on the complete experimental results, the following recommendations are made:

\begin{enumerate}
    \item \textbf{Best overall model:} Aya Expanse 8B \cite{aya_2024}. It achieved the highest dense NDCG@10 (0.6164, +23.5\%) and the second-highest BM25 NDCG@10 (0.5046, +9.2\%), making it the only model with strong improvements on \emph{both} retrieval paradigms.

    \item \textbf{Best for BM25:} Jais-2 8B \cite{sengupta_2023_jais}. It achieved the highest BM25 NDCG@10 (0.5122, +10.8\%), uniquely benefiting from its Arabic-centric vocabulary for term-based matching.

    \item \textbf{Best for resource-constrained environments:} Qwen3-4B \cite{qwen3_2025}. At 4B parameters with no quantisation required, it achieved +14.0\% dense NDCG@10---nearly double the improvement of the 3B models---while fitting on a T4 GPU in FP16.

    \item \textbf{Temperature 0.1} was found to be optimal for deterministic pseudo-document generation, yielding +2.5\% improvement over temperature 0.7 in controlled testing (Section~\ref{sec:meth_temperature}).
\end{enumerate}

\subsection{Summary of All Experiments}
\label{sec:res_summary}

Table~\ref{tab:full_summary} provides a consolidated view of all experiments conducted in this research.

\begin{table}[htbp]
    \centering
    \caption{Summary of all experiments. Exp.~001 and 002 are baselines; Exp.~003--009 use Query2Doc with the indicated model. Bold values indicate the best result for each metric across all experiments.}
    \label{tab:full_summary}
    \begin{tabular}{@{}cllcccc@{}}
        \toprule
        \textbf{Exp.} & \textbf{Model} & \textbf{Retriever} & \textbf{NDCG@10} & \textbf{Recall@10} & \textbf{MRR} & \textbf{Status} \\
        \midrule
        001 & None (baseline) & Dense & 0.4993 & 0.6156 & 0.5328 & Baseline \\
        002 & None (baseline) & BM25 & 0.4621 & 0.5964 & 0.4836 & Baseline \\
        003 & Qwen 2.5 3B & Dense & 0.5435 & 0.6608 & 0.5742 & OK \\
        004 & Qwen 2.5 3B & BM25 & 0.4090 & 0.5384 & 0.4342 & Degraded \\
        005 & Falcon-H1 3B & Dense & 0.5359 & 0.6484 & 0.5681 & OK \\
        006 & Jais-2 8B & Dense & 0.6018 & 0.7161 & 0.6356 & Good \\
        006 & Jais-2 8B & BM25 & \textbf{0.5122} & \textbf{0.6448} & \textbf{0.5397} & \textbf{Best BM25} \\
        007 & Qwen3-4B & Dense & 0.5691 & 0.6824 & 0.6015 & Good \\
        008 & ALLaM 7B & Dense & 0.2550 & 0.3335 & 0.2708 & Dropped \\
        009 & GPT-OSS 20B & --- & --- & --- & --- & Dropped \\
        --- & Aya 8B & Dense & \textbf{0.6164} & \textbf{0.7256} & \textbf{0.6493} & \textbf{Best Dense} \\
        --- & Aya 8B & BM25 & 0.5046 & 0.6284 & 0.5377 & Good \\
        \bottomrule
    \end{tabular}
\end{table}

The experimental results demonstrate that LLM-based query enhancement via Query2Doc can improve Arabic dense retrieval by up to +23.5\% NDCG@10 using open-source models that run on consumer-grade GPUs. The technique is particularly effective for dense retrieval, while BM25 requires careful model selection and potentially retriever-specific strategies to avoid term dilution. These findings establish a strong foundation for future work on Arabic query enhancement in RAG systems.
